\SourceLeaf{282}%
minimal werden soll. Nach den Regeln der Variationsrechnung ergibt sich dafür die Gleichung
\[
\frac{\mathrm{d}n}{\mathrm{d}u}\mathord{.}\BareRoot{(1+\frac{\mathrm{d}u^2}{\mathrm{d}x^2})}\mathrm{d}x=\mathrm{d}\frac{\frac{n\mathrm{d}u}{\mathrm{d}x}}{\BareRoot{(1+\frac{\mathrm{d}u^2}{\mathrm{d}x^2})}}
\]
oder
\[
\frac{\mathrm{d}n}{\mathrm{d}u}\mathord{.}\frac{\mathrm{d}x}{\cos\psi}=\mathrm{d}\mathord{.}n\sin\psi
\]
Mit \(\tfrac{\mathrm{d}n}{\mathrm{d}u}\) ist die partielle Ableitung gemeint. Diese Formel ist exakt und gilt allgemein. Für unseren Zweck kommt jedoch nur das zwischen \(F\) und \(G\) liegende Stück der Kurve \(M\) in Betracht, an dessen sämtlichen Punkten \(u\) und \(\psi\) nur sehr kleine Werte haben können. Wir dürfen also \(1\) anstelle von \(\cos\psi\) und \(\operatorname{tang}\psi\) anstelle von \(\sin\psi\) schreiben und somit
\[
\frac{\mathrm{d}n}{\mathrm{d}u}\mathord{.}\mathrm{d}x=\mathrm{d}\mathord{.}n\operatorname{tang}\psi
\]
oder
\[
n\operatorname{tang}\psi=\int\frac{\mathrm{d}n}{\mathrm{d}u}\mathrm{d}x+\text{Const.}
\]
setzen. Zugleich dürfen wir in dieser Formel anstelle der Werte, die \(n\) und \(\tfrac{\mathrm{d}n}{\mathrm{d}u}\) auf der Linie \(M\) haben, diejenigen an den entsprechenden Punkten der Linie \(N\) (für \(u=0\) oder \(y=0\)) verwenden. Diese stimmen also mit den Werten von \(\tfrac{1}{m}\) und \(-\tfrac{\mathrm{d}m}{mm\mathrm{d}u}=-\tfrac{\mathrm{d}m}{mm\mathrm{d}y}\) überein.

Um die weiteren Entwicklungen zu erleichtern, sollen die Abszissen jetzt vom Punkt \(F\) an gezählt werden. In diesem Punkt setzen wir also \(x=0\), in \(G\) dagegen \(x=h\). Ferner setze ich \(\tfrac{\mathrm{d}m}{m\mathrm{d}y}=l\). Dies ist im Allgemeinen eine Funktion von \(x\) und \(y\), wird hier jedoch nur mit dem auf der Linie \(N\) oder für \(y=0\) geltenden Wert betrachtet, also als Funktion von \(x\) allein. Schließlich seien \(\psi^0,m^0,l^0\) die bestimmten Werte von \(\psi,m,l\) im Punkt \(F\) und \(\psi',m',l'\) die im Punkt \(G\). Die obige Formel lautet danach
\[
\operatorname{tang}\psi=\frac{m\operatorname{tang}\psi^0}{m^0}-m\int\frac{l}{m}\mathrm{d}x
\]
wobei die Integration bei \(x=0\) beginnt. Nehmen wir nun an, dass \(l\) und \(m\) in die folgenden Reihen nach Potenzen von \(x\)
